\SourceLeaf{451}
\par\Needspace{8\baselineskip}\bigskip
\pdfbookmark[3]{Exposition, 4}{pallas-II-article4}
\begin{center}4.\end{center}
\noindent\hspace*{\parindent}It remains to develop the relation between our constants and the elements astronomers customarily use. For this purpose, let us refer the following nine directions to a spherical surface of arbitrary radius whose center coincides with the center of the Sun. These are the three coordinate axes, the line drawn from the Sun toward the planet, a line perpendicular to it in the plane of the orbit, the line perpendicular to that plane, the line of apsides directed toward perihelion, the perpendicular to it in the plane of the orbit, and the intersection of the plane of \(x,y\) with the plane of the orbit. We shall denote the points of the spherical surface corresponding to these directions by \(X,Y,Z,L,\Lambda,M,P,\Pi,N\). Thus \(X,Y,Z\) will be the poles of the great circles that represent the three fundamental planes, and we suppose that these poles are taken on the side where the coordinates are considered positive. Let \(L\) be the heliocentric position of the planet, \(\Lambda\) the point on the great circle of the orbit that is ahead of \(L\) by \(90^{0}\), \(P\) the perihelion, and \(\Pi\) the point to which \(\Lambda\) moves when \(L\) becomes \(P\). Let \(N\) be the ascending node of the orbit on the great circle whose pole is \(Z\). Finally, \(M\) will be one of the poles of the orbit. To avoid ambiguity, we suppose that the points \(P,\Pi,M\) are arranged on the sphere in the same order as the points \(X,Y,Z\). We see that for this purpose \(M\) must be chosen on the same side of the great circle \(XY\) as \(Z\), or on the opposite side, according as the heliocentric motion of the planet projected onto the great circle \(XY\) goes in the direction from \(X\) toward \(Y\) or in the opposite direction. We call the motion \textit{direct} in the first case and \textit{retrograde} in the second. To distinguish these two cases analytically, let \(\beta\) be the distance of \(L\) from the great circle \(XY\), or more precisely \(\beta=90^{0}-LZ\), and let \(\lambda\) be the longitude of the projection of \(L\) onto this great circle, measured from \(X\) in the direction from \(X\) toward \(Y\). We have
\[
\begin{aligned}
x&=r\cos\beta\cos\lambda\\
y&=r\cos\beta\sin\lambda\\
z&=r\sin\beta.
\end{aligned}
\]
\par From this we obtain \(x\mathrm{d}y-y\mathrm{d}x=rr\cos\beta^2\mathbin{.}\mathrm{d}\lambda\). Thus in the case of \SourceLeaf{452}direct motion \(C\) will be positive, and in the case of retrograde motion it will be negative. Hence \(C\) and \(\cos ZM\) will always have the same sign. Now we have
\[
\begin{array}{rl}
27)\quad&x=r\cos XL\\[4pt]
28)\quad&y=r\cos YL\\[4pt]
29)\quad&z=r\cos ZL
\end{array}
\]
and by a known theorem
\[
\begin{aligned}
&\cos XL\mathbin{.}\cos XM+\cos YL\mathbin{.}\cos YM+\cos ZL\mathbin{.}\cos ZM\\
&\qquad=\cos LM=0.
\end{aligned}
\]
\par Consequently we have
\[
x\cos XM+y\cos YM+z\cos ZM=0
\]
and since this equation must evidently be identical to equation 7, the quantities \(A,B,C\) must be proportional to \(\cos XM,\cos YM,\cos ZM\). By another known theorem we have \(\cos XM^2+\cos YM^2+\cos ZM^2=1\). Now we have just shown that \(\cos ZM\) and \(C\) necessarily have the same sign, and in addition \(AA+BB+CC=mp\). We must therefore conclude
\[
\begin{array}{rl}
30)\quad&A=\cos XM\mathbin{.}\BareRoot{mp}\\[4pt]
31)\quad&B=\cos YM\mathbin{.}\BareRoot{mp}\\[4pt]
32)\quad&C=\cos ZM\mathbin{.}\BareRoot{mp}.
\end{array}
\]
